% =====================================================================
%  APPENDIX A: CALCULATION TRACEABILITY  (condensed, October 2026)
%  Values as generated by make_appendix_traceability.py and
%  review_checks.py from results/processed/*.json [ref32]. The complete
%  results matrices of all 120 cases are in docs/full_results_matrices.md
%  of the repository.
% =====================================================================

\chapter{Calculation Traceability}
\label{app:traceability}

This appendix documents the input values and the calculation chain from the mission requirement to the structural margins, with numerical values for the E190 and the script that performs each step \cite{ref32}. The same relations are applied to the A320 and the A350-900 with their respective inputs. The complete result matrices of all 120 cases, the detailed derivations and the scripts are available in the repository; the file docs/full\_results\_matrices.md lists every case individually. In the method column, \emph{direct} denotes evaluation of a closed-form expression, \emph{iterative} the numerical solution of an implicit equation, and FEM a quantity obtained from the Abaqus solution.

\begin{table}[!htbp]
\centering
\caption[Input values and their sources]{Input values and their sources.}
\label{tab:app_inputs}
\footnotesize
\setlength{\tabcolsep}{4pt}
\begin{tabularx}{\textwidth}{@{}llrl>{\raggedright\arraybackslash}X@{}}
\toprule
Symbol & Quantity & Value & Unit & Source \\
\midrule
$E$, $\nu$ & Elastic constants, aluminium alloy 2219-T87 & 73.1, 0.33 & GPa, - & MIL-HDBK-5J \cite{ref30} \\
$\alpha$ & Thermal expansion coefficient & $22.3\times10^{-6}$ & 1/K & MIL-HDBK-5J \cite{ref30}, room temperature \\
$\rho$ & Density & 2850 & kg/m$^3$ & MIL-HDBK-5J \cite{ref30} \\
$\kappa$ & Thermal conductivity & 116 & W/(m\,K) & MIL-HDBK-5J \cite{ref30}, room temperature \\
$F_{\mathrm{ty}}$ & Tensile yield strength, room temperature & 393 & MPa & MIL-HDBK-5J \cite{ref30} \\
- & Design margin on both allowables & 1.5 & - & Project margin, Section~\ref{sec:failure_criteria} \\
$\rho_{\mathrm{LH2}}$ & Liquid hydrogen density & 70.8 & kg/m$^3$ & Leachman et al.\ \cite{ref14} \\
$p$, $p_{\mathrm{min}}$ & Design pressure differences, strength and stability & 1.5, 0.187 & bar & Sections~\ref{ssec:boiloff_requirement} and~\ref{sec:failure_criteria} \\
- & Fill pressure and venting pressure, absolute & 1.2, 1.448 & bar & Winnefeld et al.\ \cite{ref12} \\
$T_{\mathrm{liq}}$, $T_{\mathrm{ref}}$ & Fluid sink temperature and stress-free reference temperature & 20.3 & K & Bagarello et al.\ \cite{ref29}; Section~\ref{sec:coupling} \\
$h_{\mathrm{L}}$, $h_{\mathrm{V}}$ & Heat transfer coefficients, liquid-contact and ullage region & 1000, 1.46 & W/(m$^2$\,K) & Bagarello et al.\ \cite{ref29} \\
$\mathrm{LHV}$ & Lower heating values of Jet A and LH2 & 43, 120 & MJ/kg & Tiwari et al.\ \cite{ref2} \\
$2R/D$ & Tank-to-fuselage diameter ratio & 0.8648 & - & Size-class scaling assumption, Section~\ref{sec:philosophy} \\
$a/b$ & End-cap aspect ratio & 1.6 & - & As Bagarello et al.\ \cite{ref29} \\
$\phi_{\mathrm{f,mission}}$ & Mission fill level, determines the volume & 0.90 & - & Section~\ref{ssec:boiloff_requirement} \\
- & Fill level used for the heat leak budget only & 0.97 & - & Minimum ullage, Winnefeld et al.\ \cite{ref12} \\
$t_{\mathrm{hold}}$ & No-vent hold time & 14.7 & h & Adopted budget, Section~\ref{ssec:boiloff_requirement} \\
$k_{\mathrm{eff}}$ & Effective thermal conductivity of the MLI & 0.0003 & W/(m\,K) & Adopted engineering value \\
$D$, $m_{\mathrm{JetA}}$ & E190 fuselage diameter and kerosene mass & 3.01, 14\,243 & m, kg & Table~\ref{tab:jeta_basis} \\
\bottomrule
\end{tabularx}
\end{table}

\begin{table}[!htbp]
\centering
\caption[Calculation chain for the E190]{Calculation chain for the E190, from the mission requirement to the structural margins. Loads and stresses are for the emergency landing at $\phi_{\mathrm{f}} = 0.90$. Scripts: final\_numbers.py, boiloff\_requirement\_derivation.py, insulation\_sizing.py, outer\_jacket\_closed\_form.py, E190\_MASTER\_\allowbreak thermo\_\allowbreak structural.py and review\_checks.py \cite{ref32}.}
\label{tab:app_chain}
\footnotesize
\setlength{\tabcolsep}{3pt}
\begin{tabularx}{\textwidth}{@{}c>{\raggedright\arraybackslash}p{3.0cm}>{\raggedright\arraybackslash}X r l@{}}
\toprule
\# & Quantity & Relation & Result & Method \\
\midrule
1 & Hydrogen mass & $m_{\mathrm{LH2}} = m_{\mathrm{JetA}}\,\mathrm{LHV}_{\mathrm{JetA}}/\mathrm{LHV}_{\mathrm{LH2}}$ & 5103.8 kg & Direct \\
2 & Radius and cap height & $R = 0.8648\,D/2$; $h_{\mathrm{cap}} = R/1.6$ & 1.3015; 0.8134 m & Direct \\
3 & Volume and barrel length & $V = m_{\mathrm{LH2}}/(0.9\,\rho_{\mathrm{LH2}}) = \pi R^2 L + \tfrac{4}{3}\pi R^2 h_{\mathrm{cap}}$ & 80.10 m$^3$; 13.9675 m & Direct \\
4 & Shell area & Equation~\eqref{eq:cap_area} plus $2\pi R L$ & 130.44 m$^2$ & Direct \\
5 & Heat leak budget & Equation~\eqref{eq:constant_volume_balance}, parahydrogen, fill level 0.97 & 702.6 W & Iterative \\
6 & Insulation heat flux & $q_{\mathrm{MLI}} = Q/A$ & 5.39 W/m$^2$ & Direct \\
7 & MLI thickness & thermal resistances in series, $\sum R_i = (293 - 20)/Q$ & 14.93 mm & Iterative \\
8 & Outer jacket, ring spacing 0.25 m & collapse pressure \cite{ref33}, factor 1.5 on \SI{1}{atm} & 3.45 mm & Iterative \\
9 & Shell and liquid mass & $\rho A t$; $\phi_{\mathrm{f}}\rho_{\mathrm{LH2}}V$ & 479; 5104 kg & Direct \\
10 & Resultant load factor and inclination & $\lvert\mathbf{n}\rvert$; $\cos\theta = n_z/\lvert\mathbf{n}\rvert$ & 11.22; \SI{57.7}{\degree} & Direct \\
11 & Applied resultant force & $(m_{\mathrm{shell}} + m_{\mathrm{liq}})\lvert\mathbf{n}\rvert g$ & 614.8 kN & Direct \\
12 & Free surface and liquid depth & plane normal to $\hat{\mathbf{g}}$ enclosing $\phi_{\mathrm{f}}V$, volume sampling & 11.45 m & Iterative \\
13 & Pressure on the wall & Equation~\eqref{eq:pressure} & up to 2.392 bar & Direct \\
14 & Temperature and stress fields & discretised energy balance, then $\mathbf{K}\mathbf{u} = \mathbf{F}$ & fields & FEM \\
15 & Far-field von Mises stress & maximum over elements more than \SI{0.5}{\metre} from the rings & 256.9 MPa & FEM \\
16 & Yield reserve factor & $(F_{\mathrm{ty}}/1.5)/\sigma_{\mathrm{vM}} = 262.0/256.9$ & 1.020 & Direct \\
17 & Critical buckling stress & Equation~\eqref{eq:sp8007}: $\sigma_{\mathrm{cl}} = 44.27$ MPa, $\varphi = 1.986$, $\gamma = 0.2226$ & 9.85 MPa & Direct \\
18 & Minimum axial membrane stress, \SI{0.187}{\bar} & far-field barrel & $-5.95$ MPa & FEM \\
19 & Buckling reserve factor & $(\sigma_{\mathrm{cr}}/1.5)/\lvert\sigma_{x,\mathrm{min}}\rvert = 6.570/5.945$ & 1.105 & Direct \\
20 & Sized wall thickness & iteration on $t$ until the lower reserve factor lies between 1.00 and 1.06 & 1.289 mm & Iterative, FEM in loop \\
\bottomrule
\end{tabularx}
\end{table}

The same chain gives critical buckling stresses of 8.80 and \SI{33.16}{\mega\pascal} for the A320 and the A350-900, and the wall thicknesses of Table~\ref{tab:sized_walls}.

The cooldown values quoted in Section~\ref{sec:coupling} are reference values, not inputs, since the stress-free state of the analysis is the cold operating state. They are based on the measured integrated thermal contraction of an aluminium alloy between \SI{293}{\kelvin} and \SI{20}{\kelvin}, \SI{0.415}{\percent} for alloy 6061 \cite{ref34}: under full constraint this strain corresponds to $E\varepsilon = \SI{304}{\mega\pascal}$ in the uniaxial and $E\varepsilon/(1-\nu) = \SI{453}{\mega\pascal}$ in the equibiaxial case, and the free axial contraction of the E190 vessel with a total length of \SI{15.594}{\metre} is \SI{64.8}{\milli\metre}. A constant room-temperature expansion coefficient over the entire temperature interval would instead give \SI{95}{\milli\metre}, \SI{46}{\percent} more, because the expansion coefficient decreases strongly at low temperature (hand\_calcs.py \cite{ref32}).

\begin{table}[!htbp]
\centering
\caption[Hand calculation against the finite element model]{Analytical solution compared with the finite element model at the sized wall thicknesses, given as analytical / model. Baseline at $\phi_{\mathrm{f}} = 0.10$ unless stated otherwise; the analytical values use no finite element input. The thermal entries use the room-temperature material properties of the model. Source: hand\_calcs.py \cite{ref32}.}
\label{tab:app_agreement}
\footnotesize
\setlength{\tabcolsep}{4pt}
\begin{tabular}{@{}lccc@{}}
\toprule
Quantity & E190 & A320 & A350-900 \\
\midrule
Barrel hoop stress $pR/t$ [MPa] & 151.50 / 151.46 & 163.01 / 162.95 & 68.13 / 68.05 \\
Mean barrel axial stress $pR/2t$ [MPa] & 75.75 / 75.75 & 81.51 / 81.51 & 34.06 / 34.07 \\
Hoop stress at the end-cap equator [MPa] & $-42.4$ / $-51.0$ & $-45.6$ / $-53.2$ & $-19.1$ / $-25.6$ \\
Equilibrium, emergency landing, 0.90 [kN] & 614.8 / 614.2 & 802.1 / 801.1 & 6354.8 / 6348.5 \\
Wall temperature difference $\Delta T_{\mathrm{w}} = q_{\mathrm{MLI}}/h_{\mathrm{V}}$ [K] & 3.689 / 3.683 & 4.651 / 4.644 & 7.376 / 7.364 \\
Contact-line heat flux $q_{\mathrm{MLI}}L_{\mathrm{c}}/t$, 0.90 [W/m$^2$] & 1338 / 1317 & 1527 / 1489 & 1274 / 1284 \\
Mean axial stress, emergency minus Baseline, 0.90 [MPa] & 21.38 / 21.04 & 17.48 / 17.08 & 22.98 / 22.67 \\
Minimum axial stress, emergency landing, \SI{0.187}{\bar} [MPa] & $-11.52$ / $-5.95$ & $+1.57$ / $+1.71$ & $-26.72$ / $-20.85$ \\
Mid-ring reaction including thermal bowing, room-temperature $\alpha$ [kN] & $-2.5$ / $-1.8$ & $-24.5$ / $-11.4$ & $-33.3$ / $-29.8$ \\
Bending decay length $\pi/\beta$ [m] & 0.101 / - & 0.127 / - & 0.297 / - \\
Characteristic conduction length $L_{\mathrm{c}}$ [m] & 0.320 / - & 0.353 / - & 0.671 / - \\
\bottomrule
\end{tabular}
\end{table}
